\section{Local matching and the global theta integral}
\label{sec:theta-match}

This is the theorem about two unitary Weil representations used in
\texttt{theta\_\allowbreak{}integrals\_\allowbreak{}agree\_\allowbreak{}and\_\allowbreak{}local\_\allowbreak{}matching\_\allowbreak{}of\_\allowbreak{}literature}.
We give its local maps, rank restrictions, integral domains and
normalizations. The final subsection records which estimates and
representation-theoretic results enter the deductions below.

\subsection{Spaces, signatures and sections}

Write \(E/F\) for the CM extension with involution \(\sigma\).
The base of the theorem has \(E/\mathbb Q\) Galois,
\(\mathcal O_E/\mathcal O_F\) unramified and a standard additive
character \(\psi\). At every inert place the residue characteristic
is odd, the residue-field size exceeds three, and \(\psi_v\) is
unramified. The base also specifies split embeddings at a finite
set \(P\) of places, unramified for \(\psi\), and at a distinguished
split place \(v_1\). The two Hermitian spaces \(V,V'\) and the
skew-Hermitian space \(W\) have dimensions \(a,a',d=a+a'\) with
\(a'<a\) and \(d\) even. For the signature data in the source the
pairs are \((a-p,p),(q,a'-q),(d-p-q,p+q)\), respectively. Its
inequalities are
\[
 2(p+q)+1<(a-a')/2,\qquad 2(p+q)+1<a',\qquad
 d<m_\alpha(w)-m_\beta(w)\quad\text{at every real place }w .
\]
A pair of Hecke characters \(\alpha,\beta\), with the prescribed
infinite exponents and restrictions to the base ideles, determines
\(\nu=\alpha/\beta\). Write \(T=S_0\cup Q\subset P\), where \(S_0,Q\)
are disjoint. The characters \(\alpha,\beta\) are unramified outside
\(T\), and \(\nu\) is ramified at places of \(E\) over \(T\).
The distinguished place \(v_1\) lies outside \(T\) and splits
completely in the sense of the base. The Weil representations
\(\omega,\omega'\) use the pairs of splitting characters
\((\beta,\beta/\alpha)\), \((\alpha,\alpha/\beta)\), respectively.
Fix the models, polarizations, Haar measures and rational fundamental
domains specified with these data.

For \(Y=V\) or \(V'\) of dimension \(r\), and for a Schwartz
function \(\Phi\) in its adelic model, put
\[
 \Theta_Y(\Phi;g,h)=
   \sum_{x\in F^{r\times d}}[\omega_Y(g,h)\Phi](x),\qquad
 I_Y(\Phi)(h)=
   \int_{U(Y)(F)\backslash U(Y)(\mathbb A_F)}
         \Theta_Y(\Phi;g,h)\,dg .
\]
The second integral here uses the model's fundamental domain and measure.
The small-side subspace \(\mathcal S^\circ\) is spanned by pure tensors
that are vacuum vectors at almost all finite places. At \(v_1\), in
the specified split realization, their local factor vanishes on all
matrices of rank less than \(a'\). At infinity they have Schwartz
factors. Denote by \(\mathcal P\) the polynomial Fock subspace
in the corresponding archimedean model.

\subsection{The common quotients at finite places}

For every finite \(v\) there is a nonzero irreducible complex space
\(\Pi_v\) with maps \(p_v,p'_v\) onto it from the respective
Schwartz--Bruhat spaces. Both maps annihilate the linear span of
differences under the first member of the dual pair. Equality of
images is preserved by the action of the second member. At \(v\notin T\)
the two vacuum vectors have equal, nonzero images. For a good split
place \(v\notin T\) in the sense of the theorem (lying over a rational
prime split completely in \(E\)), the kernel on either Schwartz space
equals this span of differences. The small-side functions supported
on the full-rank locus in split coordinates surject onto \(\Pi_v\).
Every linear functional on the small Schwartz space arising from a
functional on the ambient functions and annihilating the differences
there factors through \(p'_v\) on Schwartz functions.

Here is the realization of the matching condition at \(v\in T\).
Let \(k=F_v\), choose the designated embedding \(E\longrightarrow k\),
and write \(\varpi\) for a uniformizer, \(\mathcal O\) for the
valuation ring, \(K_r=\operatorname{GL}_r(\mathcal O)\).
Use Haar measures of mass one on \(K_r\).
Use \(\alpha,\beta,\nu\) also for the local characters in these
coordinates. Put \(X_0=[1_a\ \ 0]\) and \(Y_0=[0\ \ 1_{a'}]\).
The argument of both transforms below is \(h\in\operatorname{GL}_d(k)\).
With \(h^{-t}=(h^{-1})^t\), they are
\begin{align*}
 P_A(\phi)(h)
   &=\beta(\det h)|\det h|^{-a/2}
     \mathop{\sum\nolimits^{\mathrm{fin}}}_{n_1\geq\cdots\geq n_a}
       \int_{K_a\operatorname{diag}(\varpi^{n_i})K_a}
       \phi(gX_0h^{-t})\nu(\det g)|\det g|^{d/2}\,dg,\\
 P_B(\phi')(h)
   &=\alpha(\det h)|\det h|^{a'/2}
       \int_{\operatorname{GL}_{a'}(k)}
       \phi'(gY_0h)\nu(\det g)|\det g|^{d/2}\,dg .
\end{align*}
The first expression uses the source's finite-support sum of shell
integrals indexed by decreasing integer tuples. The ramified
calculation establishes finite support on the Schwartz functions
to which the linearity argument is applied. The local Weil coordinate
changes \(\iota,\iota'\) send the two vacuum functions to indicators
of integral matrices of their respective sizes. The matching criterion is
\[
       p_v\Phi=p'_v\Phi'
        \quad\Longleftrightarrow\quad
       P_A(\iota\Phi)=P_B(\iota'\Phi')
           \quad\text{as functions on }\operatorname{GL}_d(k).
\]
There is no equality of vacuum images required at \(v\in T\).

To see how these transforms give the claimed quotient, use the
local change-of-variables and shell lemmas: both transforms are
linear on Schwartz--Bruhat functions, take values in the same
normalized parabolically induced representation, annihilate the
first group's differences, and intertwine the second group's action.
Each image contains a nonzero vector. For instance take a small
principal congruence subgroup on which \(\nu\circ\det\) is trivial.
At \(h=1\) localize the test function around the image of the
identity matrix under \(g\mapsto gX_0\) or \(g\mapsto gY_0\).
These inclusions have coordinate projections for left inverses,
so a compact open neighbourhood may be chosen with pullback
precisely the chosen subgroup. The support is then in the unit
shell, all factors in the integrand there are one if the
characteristic function is used, and its positive Haar volume is
the resulting value. This is the bump-function argument in
\texttt{p\allowbreak{}Ram\_\allowbreak{}ne\_\allowbreak{}zero} and its small-side counterpart.

The needed Bernstein irreducibility statement is the irreducibility
of the normalized product of two determinant representations of
\(\operatorname{GL}_a(k)\) and \(\operatorname{GL}_{a'}(k)\) for
smooth unitary characters when the dimensions or the characters
differ. Here \(a\ne a'\). An invariant subspace containing a
nonzero vector in this representation is the entire representation.
Both maps are therefore onto. Transport them through the split
Weil coordinates. Equality in their codomain, which is a subspace
of functions, is equivalent to pointwise equality of the transforms;
this proves the displayed ramified matching criterion.
In rank descent the source separately uses ramified multiplicity one:
for a nonzero equivariant map between the successive local models,
every other equivariant map of the same type is its scalar multiple.
It derives both local results using the geometric-lemma filtration.

For \(v\notin T\), use the split unramified maps when an embedding
\(E\longrightarrow F_v\) exists, and the inert spherical maps
otherwise. The inert input identifies the vacuum vectors in a common
spherical quotient. The source obtains it from its inert vacuum
comparison, spherical-function result and local representation
existence. These three cases exhaust the finite places. Choose the
pair of maps constructed at each \(v\); all the requirements on
them above are placewise. At infinity the two polynomial modules
surject onto a common nonzero irreducible \((\mathfrak g,K)\)-module
\(\Pi_\infty\). The maps kill first-group Lie derivatives and identify
first-group maximal-compact translates. They intertwine the
second group's Lie and maximal-compact actions on matching pairs.
The formal proof supplies this archimedean statement by its Fock model
and local lift comparison.

\subsection{The constant and the rank induction}

Two finite sums of pure tensors with polynomial archimedean factors
are polynomially matched when their images in
\(\Pi_\infty\otimes\bigotimes_{v\in U}\Pi_v\) are equal for a
finite set \(U\supset T\) outside which all their factors are
vacua. General matched data are finite sums with fixed finite
Schwartz factors whose archimedean factors are limits of sequences
of such polynomially matched sums. The small member and every one
of its approximants must lie in \(\mathcal S^\circ\).
For each choice of matching maps there is a unique \(c\ne0\) for which
\begin{equation}
       I_V(\Phi)(h)=c I_{V'}(\Phi')(h)
       \qquad\text{for all matched pairs and every }h.
       \label{eq:theta-identity}
\end{equation}

Here are the two rank changes controlling this constant. They also
specify which scalar factors may depend on local characters. Put
\(\delta=(a-a')/2\), and let \(\mathfrak b,\mathfrak c\) be the
fractional ideals of the two adapted planes for a step. For its
universal signs \(e_1,e_2,s\in\{1,-1\}\), set
\begin{align*}
 D(z)&=\nu(\mathfrak b)^{e_1}\nu(\mathfrak c)^{e_2}
       \prod_{v\in T}\nu_v(s)\epsilon_v(z),\\
 \epsilon_v(z)&=q_v^{-f_vz}\nu_v(\varpi_v)^{f_v}
      \sum_{y\in(\mathcal O_v/\varpi_v^{f_v})^\times}
        \nu_v(\widetilde y)^{-1}\psi_v(\widetilde y/\varpi_v^{f_v}).
\end{align*}
Here \(f_v\) is the conductor exponent, \(q_v\) the residue cardinality
and the sum uses lifts of the unit classes. The ideal character
values have the source's uniformizer conventions. At a beta step,
\(q>0\), the ranks change from \((a,a')\) to \((a,a'-2)\)
and \(q\) drops by one. With \(c_1\) the lower-rank constant,
\[
        c\,L^{S_0}(\nu,\delta+1)
          =[C D(-\delta)]^{-1}c_1 .
\]
At an alpha step, \(q=0,\ p>0\), the ranks change to
\((a-2,a')\), and \(p\) drops by one. In that case
\[
        c=C D(1-\delta)L^{S_0}(\nu,\delta)c_1 .
\]
Here \(L^{S_0}\) is the partial character \(L\)-function with the
specified \(S_0\) Euler factors omitted. Each step has its own
universal triple of signs. Its \(C\ne0\) is chosen after the planes,
geometry and archimedean maps, before the characters and finite
matching data.

To prove the deduction from the rank-step inputs, induct on \(p+q\).
The \(p=q=0\) case is the anisotropic polynomial identity. If \(q>0\)
use the beta step; if \(q=0\) and \(p>0\) use the alpha step.
The source's existence theorem for adapted planes supplies a
lower-rank admissible triple. Local matching supplies its models,
and the induction hypothesis supplies \(c_1\).
The comparisons of the constant terms for the zero orbit and
the relevant nonzero cone give the displayed step relations and
the polynomially matched identity for the larger pair. The
nonvanishing factor result gives \(D(z)\ne0\) at the specified
arguments and nonvanishing of the partial \(L\)-value.
Hence the new scalar is nonzero.

There is also a small-side test vector with nonzero integral.
When \(q=0\) this is the anisotropic nonvanishing conclusion.
For \(q>0\), the nonzero term in beta descent supplies a big-side
vector with nonzero integral. Choose its polynomially matched
small-side partner. Its integral must be nonzero by the polynomial
identity. Evaluating at an \(h\) where this integral is nonzero
forces any two admissible constants to be equal. To extend to
general matched data, pass to the archimedean approximants in
the definition above. The big-side integral is continuous on
each fixed finite Schwartz factor by its integrable majorant;
the small-side one is continuous along the indicated sequences
in \(\mathcal S^\circ\). There are only finitely many fixed
factors. Taking both limits in the polynomial identity proves
\eqref{eq:theta-identity}.

\subsection{Indefinite convergence and regularized representatives}

The large-side integral converges absolutely on every adelic
Schwartz function in the stated range. On the small side take
\(\Phi'\in\mathcal S^\circ\). Put \(H_Q(g)=\inf_{\gamma\in U(V')(F)}
H(\gamma g)\), where \(H\geq1\) is the matrix height of the model.
For each integer \(N\geq0\) there are \(C,k\) such that
\[
 |\Theta_{V'}(\Phi';g,h)|H_Q(g)^N\leq C H(h)^k .
\]
Moreover \(\int|\Theta_{V'}(\Phi';g,h)|\,dg\) is finite and bounded
by a fixed power of \(H(h)\) times a constant. Here the input from
unitary reduction theory consists of an integrable majorant for
each pure tensor in the full-rank subspace and a seminorm bound
with each power \(N\). The finite-sum deduction is explicit.
Write \(\Phi'=\sum_{i=1}^m f_i\otimes\phi_i\) in the defining span.
For its \(i\)-th summand use the integral estimate
\(\int|\Theta_i(g,h)|\,dg\leq C_i H(h)^{k_i}\) supplied by that
majorant. Form the finite sum, take \(k=\max_i k_i\), and use
\(\int|\sum_i\Theta_i|\leq\sum_i\int|\Theta_i|\).
The zero sum is immediate. The same argument applied pointwise to
the seminorm bounds gives the displayed rapid decrease. This also
shows directly why the integral is a moderate function of \(h\).

The subsequent twisted square comparison gives a regularized
representative for each small-side algebraic datum \(\Phi'\).
In its model let \(\mathcal R\) be the span of
\(\omega'(u,1)\Psi-\Psi\), with \(u\) in the allowed finite-adelic
subgroup and \(\Psi\) algebraic. There is \(\Phi_0\) in the same
algebraic space with \(\Phi'-\Phi_0\in\mathcal R\) and with a
convergent moderate theta integral. For any matched big-side
\(\Phi\) with \(p\Phi=p'\Phi_0\), the twisted analogue of
\eqref{eq:theta-identity} holds for \(\Phi,\Phi_0\).
The following argument from the archive shows precisely when the
integral of \(\Phi'\) itself can take its place.

Suppose \(\Psi\in\mathcal R\) and its kernel is integrable as a
function of \(g\) for every \(h\). Write
\(\Psi=\sum_i b_i(\omega'(u_i,1)\Psi_i-\Psi_i)\).
Choose a regularizer \(Z=\sum_k z_k\omega'(k,1)\), a finite sum in
the allowed subgroup with \(\deg Z=\sum_kz_k\ne0\), for this finite
list. The source constructs it at a completely split auxiliary place
outside the exceptional sets for both the kernels and the finite
translates. Its regularized vectors \(Z\Psi_i\) have convergent
moderate integrals, and for each \(i\) there is \(u'_i\) in the
same subgroup satisfying
\(Z\omega'(u_i,1)\Psi_i=\omega'(u'_i,1)Z\Psi_i\).
Right translation on the adelic quotient preserves the integral
of each integrable kernel. Thus
\[
 I_{V'}(Z\Psi)=
   \sum_i b_i\{I_{V'}(\omega'(u'_i,1)Z\Psi_i)-I_{V'}(Z\Psi_i)\}=0 .
\]
Since the kernel of \(\Psi\) itself is integrable, right translation
and the finite expansion of \(Z\) also give
\(I_{V'}(Z\Psi)=(\deg Z)I_{V'}(\Psi)\).
Division by the nonzero degree proves \(I_{V'}(\Psi)=0\).
If the original \(\Phi'\) has integrable kernel, apply this to
\(\Psi=\Phi'-\Phi_0\). Since the matching maps annihilate
\(\mathcal R\), \(p'\Phi'=p'\Phi_0\); since
\(I_{V'}(\Phi')=I_{V'}(\Phi_0)\), the scalar identity holds with
the ordinary integral of \(\Phi'\). For a general algebraic
\(\Phi'\) the comparison uses the regularized representative.

\paragraph{Inputs in the release.}
At the top edge \texttt{theorem\allowbreak{}K\_\allowbreak{}head2} takes two propositions.
\texttt{Eis\allowbreak{}BL} gives the half-plane of convergence of a doubled
Eisenstein family and weak holomorphic continuation after a finite
local pole-clearing factor. \texttt{Siegel\allowbreak{}Weil\allowbreak{}Std} gives, for a
definite first member of dimension \(m\) and second member of
dimension \(n\), the holomorphic value of the specified Eisenstein
section at \((m-n)/2\). It is the average of the theta sum over the
first group multiplied by \(1\) if \(n<m\) and \(2\) otherwise.
These two propositions are supplied by \texttt{Rg\allowbreak{}Eis.\allowbreak{}Hd.\allowbreak{}eis\allowbreak{}BL}
and \texttt{Siegel\allowbreak{}Weil\allowbreak{}Formula.\allowbreak{}Fin.\allowbreak{}siegel\allowbreak{}Weil\allowbreak{}Std}. Neither supplier
uses the nominal literature-bundle parameter in its proof.

The fourteen arguments in the source's next assembly
\texttt{Tr.\allowbreak{}close\allowbreak{}Sph} are the finite coinvariant dictionary,
the full-rank small-space coinvariant theorem, ramified
multiplicity one, the ramified finite matching maps, split
unramified finite matching maps, the spherical triple comparison,
inert finite matching maps, the second-group real Lie derivative
comparison, archimedean Lie equivariance, the common real quotient,
the unitary Siegel-set theorem, the Eisenstein continuation theorem,
the definite Siegel--Weil value, and existence of the requisite
Weil representations. The spherical input uses
\texttt{mod\allowbreak{}St2}; the ramified and split finite arguments use the
geometric-lemma filtration; the real quotient uses the local lift
and infinitesimal action. Inside \texttt{close\allowbreak{}Sph}, the
nontrivial analytic and algebraic steps cited above include
\texttt{pol\allowbreak{}Zero}, \texttt{non\allowbreak{}Zero\allowbreak{}An}, \texttt{end\allowbreak{}AB\_\allowbreak{}sph},
\texttt{loc\allowbreak{}A\_\allowbreak{}of}, \texttt{loc\allowbreak{}B\_\allowbreak{}of}, the plane towers,
\texttt{dom\allowbreak{}Weil\allowbreak{}U}, \texttt{so\allowbreak{}Semi\allowbreak{}U} and the unitary Siegel-set
estimates. Their underlying proofs are imported in this chapter.
The correspondence with the definitions is
\texttt{Weil\allowbreak{}Tensor\allowbreak{}Models.\allowbreak{}Is\allowbreak{}Matched}, \texttt{Decay},
\texttt{One\allowbreak{}Const}, \texttt{Coinv}, \texttt{Step\allowbreak{}Beta} and
\texttt{Step\allowbreak{}Alpha}.

The additional application to regularized representatives is
\texttt{twisted\_\allowbreak{}square\_\allowbreak{}theta\_\allowbreak{}integrals\_\allowbreak{}eq}; its inputs are
the two tensor-model theorems and
\texttt{Global\allowbreak{}Theory.\allowbreak{}Fun\allowbreak{}Pr.\allowbreak{}Theta\allowbreak{}Local\allowbreak{}Inputs}. The removal of the
regularizer for an already integrable kernel uses
\texttt{twist\_\allowbreak{}regulariser\_\allowbreak{}is\allowbreak{}Good},
\texttt{theta\allowbreak{}Integral\_\allowbreak{}eq\_\allowbreak{}zero\_\allowbreak{}of\_\allowbreak{}finite\allowbreak{}Relation} and
\texttt{theta\allowbreak{}Integrals\_\allowbreak{}agree\_\allowbreak{}of\_\allowbreak{}integrable\_\allowbreak{}kernel}.
The auxiliary-place construction uses placewise seesaw models,
the full-rank regularization theorem, commutation of separated
finite-place operators and existence of a completely split place
outside a prescribed finite set. These are source results used by
the argument in this subsection.
